% ══════════════════════════════════════════════════════════════════════════════
% PREÁMBULO OFICIAL — Emiliano Miranda González / ITAM Ciencia Política
% Versión 1.0 — Marzo 2026
% ESTE ARCHIVO ES INMUTABLE. No modificar colores, fuentes, ni templates.
% ══════════════════════════════════════════════════════════════════════════════

\documentclass[
  aspectratio=43,        % Relación 4:3 — estándar académico formal
  t,                     % Alineación vertical: top (no centrado)
  10pt                   % Tamaño base de fuente
]{beamer}

% ── TIPOGRAFÍA ────────────────────────────────────────────────────────────────
\usepackage{lmodern}           % Latin Modern Roman — fuente principal, serif
\usepackage[T1]{fontenc}       % Codificación de salida T1 (acentos correctos)
\usepackage[utf8]{inputenc}    % Codificación de entrada UTF-8
\usepackage{microtype}         % Mejoras de kerning y justificación

% ── MATEMÁTICAS ──────────────────────────────────────────────────────────────
\usepackage{amsmath}           % Entornos matemáticos avanzados
\usepackage{amssymb}           % Símbolos matemáticos AMS
\usepackage{bm}                % Negritas en modo matemático (\bm{x})
\usepackage{mathtools}         % Extensión de amsmath

% ── GRÁFICAS Y FIGURAS ───────────────────────────────────────────────────────
\usepackage{graphicx}          % Inclusión de imágenes externas
\usepackage{booktabs}          % Tablas académicas (\toprule, \midrule, \bottomrule)
\usepackage{tabularx}          % Tablas con columnas de ancho variable
\usepackage{multirow}          % Celdas que abarcan múltiples filas
\usepackage{array}             % Mejoras en columnas de tablas

% ── COLOR Y CAJAS ────────────────────────────────────────────────────────────
\usepackage{xcolor}            % Sistema de colores extendido
\usepackage{tcolorbox}         % Cajas de colores personalizadas
\tcbuselibrary{skins, breakable}

% ── TIKZ (para portada con rectángulo de color) ───────────────────────────────
\usepackage{tikz}

% ── LAYOUT Y UTILIDADES ──────────────────────────────────────────────────────
\usepackage{multicol}          % Múltiples columnas
\usepackage{appendixnumberbeamer} % Reinicia numeración en apéndice
% NOTA: enumitem NO se carga — conflicto con beamer. Usar \setlength{\itemsep}
% directamente o los templates de itemize de beamer.

% ── HIPERENLACES (siempre al final del preámbulo) ────────────────────────────
\usepackage{url}               % Formateo de URLs
\usepackage{hyperref}          % Hipervínculos — SIEMPRE el último paquete

% ══════════════════════════════════════════════════════════════════════════════
% PALETA CROMÁTICA — NO MODIFICAR
% ══════════════════════════════════════════════════════════════════════════════
\definecolor{ITAMgreen}{HTML}{00563F}      % Sacramento State Green — primario
\definecolor{ITAMmaroon}{HTML}{7B1113}     % UP Maroon — secundario/acento
\definecolor{ITAMgray}{HTML}{4A4A4A}       % Gris oscuro — texto secundario
\definecolor{ITAMlightgray}{HTML}{F2F2F2}  % Gris claro — fondos
\definecolor{ITAMwhite}{HTML}{FFFFFF}      % Blanco puro
\definecolor{blockbodybg}{HTML}{E8F0EC}    % Verde muy claro — fondo bloque primario
\definecolor{examplebodybg}{HTML}{F5E8E8}  % Rojo muy claro — fondo bloque secundario

% ══════════════════════════════════════════════════════════════════════════════
% TEMA BEAMER — NO MODIFICAR
% ══════════════════════════════════════════════════════════════════════════════
\usetheme{default}
\usecolortheme{default}
\usefonttheme{serif}           % Latin Modern Roman en TODO el documento
\usefonttheme{professionalfonts}
\setbeamertemplate{navigation symbols}{}   % Eliminar botones de navegación

% ── Colores de elementos ──────────────────────────────────────────────────────
\setbeamercolor{title}{fg=ITAMwhite}
\setbeamercolor{subtitle}{fg=ITAMwhite}
\setbeamercolor{author}{fg=ITAMwhite}
\setbeamercolor{institute}{fg=ITAMwhite}
\setbeamercolor{date}{fg=ITAMwhite}

\setbeamercolor{frametitle}{fg=ITAMgreen, bg=ITAMwhite}
\setbeamercolor{framesubtitle}{fg=ITAMgray}
\setbeamercolor{normal text}{fg=black, bg=ITAMwhite}
\setbeamercolor{structure}{fg=ITAMgreen}

\setbeamercolor{item}{fg=ITAMgreen}
\setbeamercolor{subitem}{fg=ITAMgray}
\setbeamercolor{subsubitem}{fg=ITAMgray}

\setbeamercolor{block title}{fg=ITAMwhite, bg=ITAMgreen}
\setbeamercolor{block body}{fg=black, bg=blockbodybg}

\setbeamercolor{block title alerted}{fg=ITAMwhite, bg=ITAMmaroon}
\setbeamercolor{block body alerted}{fg=black, bg=examplebodybg}

\setbeamercolor{block title example}{fg=ITAMwhite, bg=ITAMmaroon}
\setbeamercolor{block body example}{fg=black, bg=examplebodybg}

\setbeamercolor{section in head/foot}{fg=ITAMgray, bg=ITAMlightgray}
\setbeamercolor{page number in head/foot}{fg=ITAMgray, bg=ITAMlightgray}
\setbeamercolor{section in head/foot shaded}{fg=ITAMgray!50, bg=ITAMlightgray}

% ── Tamaños de fuente ─────────────────────────────────────────────────────────
\setbeamerfont{title}{size=\Large, series=\bfseries}
\setbeamerfont{subtitle}{size=\normalsize, series=\normalfont}
\setbeamerfont{author}{size=\normalsize}
\setbeamerfont{institute}{size=\small}
\setbeamerfont{date}{size=\small}
\setbeamerfont{frametitle}{size=\large, series=\normalfont}
\setbeamerfont{framesubtitle}{size=\small}
\setbeamerfont{block title}{size=\normalsize, series=\bfseries}
\setbeamerfont{block body}{size=\small}
\setbeamerfont{footline}{size=\tiny}

% ── Template de título de frame con línea separadora ─────────────────────────
\setbeamertemplate{frametitle}{%
  \vspace{0.4em}
  \textcolor{ITAMgreen}{\insertframetitle}%
  \ifx\insertframesubtitle\@empty\else%
    \\[0.1em]{\usebeamerfont{framesubtitle}\usebeamercolor[fg]{framesubtitle}\insertframesubtitle}%
  \fi
  \vspace{0.1em}
  \textcolor{ITAMgreen}{\hrule height 0.4pt}%
  \vspace{0.3em}
}

% ── Bullets ───────────────────────────────────────────────────────────────────
\setbeamertemplate{itemize item}{\textcolor{ITAMgreen}{\textbullet}}
\setbeamertemplate{itemize subitem}{\textcolor{ITAMgray}{\small\textbullet}}
\setbeamertemplate{itemize subsubitem}{\textcolor{ITAMgray}{\tiny\textbullet}}
\setbeamertemplate{enumerate item}{\textcolor{ITAMgreen}{\insertenumlabel.}}
\setbeamertemplate{enumerate subitem}{\textcolor{ITAMgray}{\insertenumlabel.}}
\setlength{\itemsep}{0.5em}

% ── Pie de página con barra de secciones (estilo Gary King) ──────────────────
\setbeamertemplate{footline}{%
  \leavevmode%
  \hbox{%
    \begin{beamercolorbox}[wd=0.75\paperwidth, ht=2.25ex, dp=1ex, leftskip=1em]{section in head/foot}%
      \usebeamerfont{footline}%
      \insertsectionnavigationhorizontal{0.75\paperwidth}{}{}%
    \end{beamercolorbox}%
    \begin{beamercolorbox}[wd=0.25\paperwidth, ht=2.25ex, dp=1ex, rightskip=1em, right]{page number in head/foot}%
      \usebeamerfont{footline}%
      \insertframenumber{} / \inserttotalframenumber%
    \end{beamercolorbox}%
  }%
  \vskip0pt%
}

% ── Portada con rectángulo de color (estilo Gary King / Lucardi) ──────────────
% Portada con posicionamiento absoluto por nodos TikZ (versión robusta,
% tomada del deck de referencia que compila correctamente).
\setbeamertemplate{title page}{%
  \begin{tikzpicture}[remember picture, overlay]
    % Mitad superior verde — bloque del título
    \fill[ITAMgreen]
      ([yshift=0pt]current page.north west) rectangle
      ([yshift=-0.5\paperheight]current page.north east);
    % Título + subtítulo
    \node[anchor=center, text width=0.86\paperwidth, align=center, text=ITAMwhite]
      at ([yshift=-0.18\paperheight]current page.north) {%
        \usebeamerfont{title}\inserttitle\par
        \vspace{0.7em}
        {\usebeamerfont{subtitle}\insertsubtitle}%
      };
    % Autor
    \node[anchor=center, text width=0.86\paperwidth, align=center, text=ITAMwhite]
      at ([yshift=-0.40\paperheight]current page.north) {%
        \usebeamerfont{author}\insertauthor%
      };
    % Mitad inferior blanca — institución y fecha
    \node[anchor=center, text width=0.86\paperwidth, align=center, text=black]
      at ([yshift=-0.75\paperheight]current page.north) {%
        \insertinstitute\par
        \vspace{0.8em}
        \insertdate%
      };
  \end{tikzpicture}%
}

% ══════════════════════════════════════════════════════════════════════════════
% METADATOS — autor e institución (un solo autor)
% ══════════════════════════════════════════════════════════════════════════════
\author[E. Miranda]{Emiliano Miranda González}
\institute[ITAM]{%
  Instituto Tecnológico Autónomo de México\\
  Departamento Académico de Ciencia Política\\[0.4em]
  {\fontfamily{cmtt}\selectfont
   \href{mailto:emiliano.miranda@itam.mx}{emiliano.miranda@itam.mx}}%
}

\title[Reelección e incumbencia]{El alcalde en el umbral}
\subtitle{(Des)ventaja del partido incumbente en los municipios mexicanos}
\date{%
  Seminario de Tópicos de Políticas Públicas\\[0.2em]
  \textit{Ciudad de México, mayo de 2026}%
}

% ══════════════════════════════════════════════════════════════════════════════
% HYPERREF
% ══════════════════════════════════════════════════════════════════════════════
\hypersetup{
  colorlinks  = true,
  linkcolor   = ITAMgreen,
  urlcolor    = ITAMgreen,
  citecolor   = ITAMgray,
  pdfauthor   = {Emiliano Miranda González},
  pdftitle    = {El alcalde en el umbral: reeleccion e incumbencia municipal en Mexico},
  pdfsubject  = {Ciencia Politica},
  pdfkeywords = {ITAM, reeleccion, incumbente, RDD, diferencias en discontinuidades, Mexico},
  pdfcreator  = {LaTeX/Beamer, Overleaf},
  unicode     = true
}

% ══════════════════════════════════════════════════════════════════════════════
\begin{document}
% ══════════════════════════════════════════════════════════════════════════════

% ── Portada ───────────────────────────────────────────────────────────────────
{
\setbeamertemplate{footline}{}
\begin{frame}[plain, noframenumbering]
  \titlepage
\end{frame}
}

% ══════════════════════════════════════════════════════════════════════════════
\section{Motivación y teoría}
% ══════════════════════════════════════════════════════════════════════════════
{
\setbeamercolor{background canvas}{bg=ITAMlightgray}
\begin{frame}[plain, noframenumbering]
  \vfill
  \begin{center}
    \textcolor{ITAMgreen}{\Large\bfseries Motivación y teoría}\\[0.5em]
    \textcolor{ITAMgray}{\normalsize De la ventaja del incumbente a su maldición}
  \end{center}
  \vfill
\end{frame}
}

% ── 1.1 ───────────────────────────────────────────────────────────────────────
\begin{frame}{La paradoja del incumbente}

  \begin{itemize}
    \item En las democracias consolidadas, gobernar confiere ventaja: Lee (2008)
          estima cerca de \textcolor{ITAMgreen}{$+8$ puntos} de voto mediante un
          \emph{close-election} RDD (EE.\,UU.)
    \item En las democracias en desarrollo el hallazgo se invierte: gobernar
          \textcolor{ITAMmaroon}{deprime} las perspectivas del partido
    \item Klašnja y Titiunik (2017) lo llaman la \emph{maldición del incumbente}:
          partidos débiles más límites de mandato anulan la rendición de cuentas
    \item En México, Lucardi y Rosas (2016) estiman una desventaja de
          \textcolor{ITAMmaroon}{23 a 29 puntos} en la probabilidad de retener
  \end{itemize}

  \vspace{0.4em}
  \begin{block}{Pregunta de investigación}
    \textit{¿Atenuó el restablecimiento de la reelección consecutiva municipal
    (reforma de 2014) la desventaja del partido incumbente en México?}
  \end{block}

\end{frame}

% ── 1.2 ───────────────────────────────────────────────────────────────────────
\begin{frame}{Mecanismos: extracción de rentas y disciplina}

  \begin{columns}[T]

    \column{0.48\textwidth}
    \textcolor{ITAMgreen}{Por qué la maldición}
    \vspace{0.3em}
    \begin{itemize}
      \item Límite de mandato $\rightarrow$ horizonte corto $\rightarrow$ alcalde
            \emph{lame duck}
      \item Extracción de rentas (Klašnja 2015); el votante castiga al partido
      \item Sanción de la corrupción revelada (Ferraz y Finan 2008)
    \end{itemize}

    \column{0.48\textwidth}
    \textcolor{ITAMmaroon}{Qué predice la reforma}
    \vspace{0.3em}
    \begin{itemize}
      \item Reelegibilidad $\rightarrow$ horizonte largo $\rightarrow$ disciplina
      \item Magar (2026): la maldición se revierte cuando el alcalde busca
            reelegirse (voto personal)
      \item Pero ese efecto confunde institución y autoselección del candidato
    \end{itemize}

  \end{columns}

  \vspace{0.6em}
  {\footnotesize \textcolor{ITAMgray}{%
    Lucardi y Rosas (2016) añaden la coordinación duvergeriana de opositores
    en torno al segundo lugar de una elección cerrada.}}

\end{frame}

% ── 1.3 ───────────────────────────────────────────────────────────────────────
\begin{frame}{La reforma de 2014 y la contrarreforma de 2030}

  \begin{itemize}
    \item Artículo 115 constitucional (2014): hasta dos periodos consecutivos
          para alcaldes, tras \textcolor{ITAMgreen}{ocho décadas} de prohibición
    \item \textcolor{ITAMgreen}{Adopción escalonada} por entidad: 22 estados en
          2018, 3 en 2019, 2 en 2021; Hidalgo y Veracruz \emph{nunca} adoptaron
    \item El objetivo causal es el efecto de la \emph{reelegibilidad como
          institución} —la disponibilidad de la regla—, que ninguna estimación
          ha aislado
    \item Una \textcolor{ITAMmaroon}{contrarreforma} restituirá la prohibición
          en 2030: saber si la regla restituyó la rendición de cuentas es
          dirimente para entender qué se cede con ese retorno
  \end{itemize}

\end{frame}

% ══════════════════════════════════════════════════════════════════════════════
\section{Datos y diseño}
% ══════════════════════════════════════════════════════════════════════════════
{
\setbeamercolor{background canvas}{bg=ITAMlightgray}
\begin{frame}[plain, noframenumbering]
  \vfill
  \begin{center}
    \textcolor{ITAMgreen}{\Large\bfseries Datos y diseño empírico}\\[0.5em]
    \textcolor{ITAMgray}{\normalsize Diferencias en discontinuidades}
  \end{center}
  \vfill
\end{frame}
}

% ── 2.1 ───────────────────────────────────────────────────────────────────────
\begin{frame}{Datos y construcción del panel}

  \begin{itemize}
    \item Base seccional de Calderón-Hernández et al.\ (2025), versión extendida
          a 2024: \textcolor{ITAMgreen}{551{,}882} observaciones de
          sección-elección, 2{,}041 municipios
    \item Calendario de adopción de Magar (2024) para la variación estado-año
    \item Colapso a municipio-elección (conteos por suma, \emph{shares}
          recalculados); panel apilado: dos filas por elección
          (\emph{ganador} y \emph{segundo lugar})
    \item Umbral $c = 0$ sobre el margen de victoria firmado; las exclusiones de
          validez dejan \textcolor{ITAMgreen}{29 entidades, 1{,}952 municipios}
  \end{itemize}

  \vspace{0.3em}
  {\footnotesize \textcolor{ITAMgray}{%
    Variable dependiente: \texttt{outcome\_rdd}, la fracción de voto válido del
    partido focal. Tratamiento: \texttt{reforma\_st}. Variable de asignación:
    \texttt{margin\_signed}.}}

\end{frame}

% ── 2.2 — estimando ───────────────────────────────────────────────────────────
\begin{frame}{El estimando: un doble diferencial}

  \begin{itemize}
    \item El RDD lee la (des)ventaja partidista como el \emph{salto} en $c=0$;
          el componente de diferencias compara \emph{saltos}, no municipios
    \item Cuatro celdas: tratados y control, antes y después de la reforma
  \end{itemize}

  \vspace{0.4em}
  \[
    \tau \;=\;
    \underbrace{\big[\,\beta^{\text{post}}_{T}-\beta^{\text{pre}}_{T}\,\big]}_{\text{1.\textsuperscript{er} diferencial (tratados)}}
    \;-\;
    \underbrace{\big[\,\beta^{\text{post}}_{C}-\beta^{\text{pre}}_{C}\,\big]}_{\text{2.\textsuperscript{o} diferencial (control)}}
  \]

  \vspace{0.2em}
  \begin{block}{Supuesto identificador (Grembi et al.\ 2016; Picchetti et al.\ 2026)}
    \small La discontinuidad confundidora es estable en el tiempo; el segundo
    diferencial neutraliza el choque común de 2018 (realineamiento de MORENA).
  \end{block}

\end{frame}

% ── 2.3 — por qué este RDD ────────────────────────────────────────────────────
\begin{frame}{¿Por qué esta variante de RDD?}

  \begin{itemize}
    \item Estimador de polinomio local (\emph{rdrobust}), orden lineal, ancho de
          banda MSE-óptimo y corrección de sesgo robusta (Calonico, Cattaneo y
          Titiunik 2014)
    \item Se descarta el polinomio global de orden alto: ruidoso y de cobertura
          deficiente (Gelman e Imbens 2019)
    \item Se descarta la aleatorización local: el margen es cuasi-continuo y
          denso en el umbral
    \item Inferencia con \textcolor{ITAMgreen}{$\approx 29$ conglomerados}
          estatales: \emph{wild cluster bootstrap} Rademacher (MacKinnon,
          Nielsen y Webb 2023)
  \end{itemize}

  \vspace{0.3em}
  {\footnotesize \textcolor{ITAMgray}{%
    Validación: densidad de Cattaneo, Jansson y Ma (2020) por partido; balance
    de covariables predeterminadas (Marshall 2024); robustez de ancho de banda.}}

\end{frame}

% ══════════════════════════════════════════════════════════════════════════════
\section{Resultados}
% ══════════════════════════════════════════════════════════════════════════════
{
\setbeamercolor{background canvas}{bg=ITAMlightgray}
\begin{frame}[plain, noframenumbering]
  \vfill
  \begin{center}
    \textcolor{ITAMgreen}{\Large\bfseries Resultados}\\[0.5em]
    \textcolor{ITAMgray}{\normalsize Un cero impreciso, y por qué}
  \end{center}
  \vfill
\end{frame}
}

% ── 3.1 — efecto agregado ─────────────────────────────────────────────────────
\begin{frame}{El efecto agregado: indistinguible de cero e impreciso}

  \begin{itemize}
    \item El diseño \textcolor{ITAMgreen}{supera el balance}: ninguna covariable
          predeterminada salta en el umbral
    \item Doble diferencial agregado: $\hat{\tau} = 0.008$ ($\approx +0.9$
          puntos), \emph{wild bootstrap} $p = 0.84$
    \item Su intervalo de confianza abarca \textcolor{ITAMmaroon}{$[-0.34,\,
          0.35]$}: $\pm 34$ puntos; sin precisión útil
  \end{itemize}

  \vspace{0.3em}
  \begin{table}\centering\small
    \begin{tabular}{lcc}
      \toprule
       & Pre-reforma & Post-reforma \\
      \midrule
      $\beta$: tratados            & $0.010\;(0.015)$ & $0.032\;(0.022)$ \\
      $\beta$: control (HGO+VER)   & $0.035\;(0.043)$ & $0.072^{***}\;(0.022)$ \\
      \midrule
      $\hat{\tau}$ doble diferencial & \multicolumn{2}{c}{$0.008\;(0.034)$} \\
      \bottomrule
    \end{tabular}
  \end{table}

  \vspace{0.2em}
  {\footnotesize \textcolor{ITAMgray}{%
    RDD local-lineal, kernel triangular, $h$ MSE-óptimo; SE entre paréntesis.
    $^{***}\,p<0.01$.}}

\end{frame}

% ── 3.2 — la causa de la imprecisión ──────────────────────────────────────────
\begin{frame}{La imprecisión es estructural: dos estados de control}

  \begin{itemize}
    \item El segundo diferencial descansa sobre los \textcolor{ITAMmaroon}{dos
          únicos} estados nunca tratados, Hidalgo y Veracruz
    \item Sus celdas post-2018 tienen muy pocas contiendas cerradas (hasta
          \textcolor{ITAMmaroon}{16} observaciones de un lado del umbral)
    \item La discontinuidad de control queda casi sin identificar; los
          estimadores por partido heredan esa imprecisión
  \end{itemize}

  \vspace{0.5em}
  \begin{alertblock}{Lectura honesta}
    \small Los $\hat{\tau}$ por partido tienen signo negativo y magnitud aparente
    grande —PAN $-0.34$, PRD $-0.18$, PRI $-0.00$—, pero ninguno es significativo
    y sus intervalos incluyen valores inverosímiles (el del PAN, hasta $-1.23$).
    \emph{No} admiten lectura causal.
  \end{alertblock}

\end{frame}

% ── 3.3 — lo informativo: la maldición por celda ──────────────────────────────
\begin{frame}{Lo informativo: la maldición sí aparece por celda}

  \begin{itemize}
    \item Antes de la reforma, los tres partidos muestran desventaja del
          incumbente en voto: PRI $-3.1$, PAN $-5.7$, PRD $-3.4$ puntos
    \item Congruente en signo con Klašnja y Titiunik (2017) y Lucardi y Rosas
          (2016)
    \item El primer diferencial intra-tratados del PAN va de
          \textcolor{ITAMgreen}{$-5.7$ a $+1.8$} puntos: atenuación compatible
          con disciplina en un partido cohesionado
    \item El PRI (partido-máquina) no se mueve; el PRD se desordena tras la
          irrupción de MORENA
  \end{itemize}

  \vspace{0.3em}
  {\footnotesize \textcolor{ITAMgray}{%
    Conjetural: el primer diferencial confunde la reforma con el realineamiento
    de 2018, y el segundo —que debía separarlos— carece de potencia.}}

\end{frame}

% ── 3.4 — tabla por partido + densidad ────────────────────────────────────────
\begin{frame}{Doble diferencial por partido y prueba de densidad}

  \begin{columns}[T]
    \column{0.56\textwidth}
    {\small Doble diferencial $\hat{\tau}$}
    \begin{table}\centering\footnotesize
      \begin{tabular}{lccc}
        \toprule
         & PRI & PAN & PRD \\
        \midrule
        Mayor único  & $-0.005$ & $-0.336$ & $-0.183$ \\
                     & $(0.048)$ & $(0.036)$ & $(0.037)$ \\
        Partido puro & $0.082$ & $-0.356^{*}$ & $-0.176$ \\
                     & $(0.073)$ & $(0.035)$ & $(0.036)$ \\
        \bottomrule
      \end{tabular}
    \end{table}

    \column{0.40\textwidth}
    {\small Densidad (McCrary)}
    \begin{table}\centering\footnotesize
      \begin{tabular}{lc}
        \toprule
        Partido & $p$ \\
        \midrule
        PRI & $0.376$ \\
        PAN & \textcolor{ITAMmaroon}{$0.010$} \\
        PRD & $0.776$ \\
        \bottomrule
      \end{tabular}
    \end{table}
  \end{columns}

  \vspace{0.4em}
  \begin{itemize}
    \item La densidad \textcolor{ITAMmaroon}{rechaza} la continuidad para el PAN:
          posible \emph{sorting}; su resultado se lee con reserva
    \item $\hat{\tau}$ es sensible al ancho de banda ($-0.013$ a $+0.062$ entre
          $0.5h$ y $2h$): no se asienta en un signo
  \end{itemize}

  \vspace{0.2em}
  {\footnotesize \textcolor{ITAMgray}{$^{*}\,p<0.10$; SE entre paréntesis;
    $p$ por \emph{wild cluster bootstrap}.}}

\end{frame}

% ── 3.5 — MORENA ──────────────────────────────────────────────────────────────
\begin{frame}{Exploración: MORENA, el partido dominante}

  \begin{itemize}
    \item Como incumbente, MORENA aparece solo desde 2018
          (\textcolor{ITAMmaroon}{1{,}117 de 1{,}186} observaciones en
          \emph{treated-post}, ninguna pre-reforma)
    \item El doble diferencial es \emph{inestimable}: un partido nacido tras la
          reforma no tiene un ``antes'' que sirva de contrafactual
    \item Lo identificable —el nivel de su discontinuidad en el régimen
          reformado— es \textcolor{ITAMmaroon}{negativo}: $-0.050$ ($p=0.035$)
          en la muestra completa; $-0.048$ en \emph{treated-post}
          (\emph{bootstrap} $p=0.083$)
    \item Ni la hegemonía nacional ni los \emph{coattails} eximen a MORENA de la
          maldición municipal —pero el resultado no es atribuible a la reforma
  \end{itemize}

\end{frame}

% ══════════════════════════════════════════════════════════════════════════════
\section{Conclusión}
% ══════════════════════════════════════════════════════════════════════════════
{
\setbeamercolor{background canvas}{bg=ITAMlightgray}
\begin{frame}[plain, noframenumbering]
  \vfill
  \begin{center}
    \textcolor{ITAMgreen}{\Large\bfseries Conclusión}\\[0.5em]
    \textcolor{ITAMgray}{\normalsize Qué permite y qué no permite concluir}
  \end{center}
  \vfill
\end{frame}
}

% ── 4.1 — conclusiones ────────────────────────────────────────────────────────
\begin{frame}{Conclusiones}

  \begin{itemize}
    \item El efecto institucional de la reelegibilidad, neto de la autoselección
          del alcalde, \textcolor{ITAMmaroon}{no se identifica con precisión
          útil} en 2018--2024
    \item No hay evidencia robusta de que la reforma restaurara la rendición de
          cuentas partidista municipal; el único indicio —el PAN— es no
          significativo e inseparable del realineamiento
    \item Diálogo con Magar (2026): donde él ve reversión \emph{cuando el alcalde
          contiende}, el efecto promedio de la \emph{mera disponibilidad} de la
          regla no desplaza la ventaja del partido
    \item Limitación dirimente: \textcolor{ITAMmaroon}{dos estados de control};
          fortalezas: balance satisfecho, \emph{wild bootstrap}, atribución de
          coaliciones depurada
  \end{itemize}

\end{frame}

% ── 4.2 — política pública ────────────────────────────────────────────────────
\begin{frame}{Implicación de política pública}

  \begin{block}{La contrarreforma de 2030}
    La restitución de la prohibición no sacrifica una ganancia de rendición de
    cuentas \emph{demostrada} —aquí no se acredita—, mas tampoco puede ampararse
    en evidencia de que la reelección fuese inocua: el diseño no descarta efectos
    sustantivos.
  \end{block}

  \vspace{0.4em}
  \begin{itemize}
    \item Como advierte Magar, la contrarreforma \emph{clausura el experimento}:
          cancela el horizonte de datos que habría permitido dirimirlo
    \item \textcolor{ITAMgreen}{Trabajo futuro}: ampliar el contrafactual con
          unidades aún-no-tratadas; usar la probabilidad de retención como
          desenlace; explorar la heterogeneidad seccional con un RDD basado en
          aleatorización
  \end{itemize}

  \vspace{0.5em}
  \centering
  {\large \textcolor{ITAMgreen}{Gracias. ¿Preguntas?}}

\end{frame}

% ══════════════════════════════════════════════════════════════════════════════
% APÉNDICE
% ══════════════════════════════════════════════════════════════════════════════
\appendix

\begin{frame}[plain, noframenumbering]
  \centering
  \vfill
  {\Large \textcolor{ITAMgreen}{\bfseries Apéndice}}
  \vfill
\end{frame}

% ── A1 — composición de la muestra ────────────────────────────────────────────
\begin{frame}[noframenumbering]{A1. Composición de la muestra por celda}

  \begin{table}\centering\small
    \begin{tabular}{lrrr}
      \toprule
      Celda & $N$ obs. & Municipios & Estados \\
      \midrule
      Tratados, pre-reforma            & 23{,}084 & 1{,}649 & 27 \\
      Tratados, post-reforma           & 7{,}990  & 1{,}579 & 26 \\
      Control (HGO+VER), pre-2018      & 3{,}658  & 296     & 2  \\
      Control (HGO+VER), desde 2018    & 746      & 293     & 2  \\
      \bottomrule
    \end{tabular}
  \end{table}

  \vspace{0.3em}
  {\footnotesize \textcolor{ITAMgray}{%
    Panel apilado (ganador + segundo lugar). Excluye CDMX, Nayarit y Tlaxcala.
    La escasez de las dos últimas filas es la limitación central del diseño.}}

\end{frame}

% ── A2 — MORENA ───────────────────────────────────────────────────────────────
\begin{frame}[noframenumbering]{A2. MORENA: discontinuidad post-reforma}

  \begin{table}\centering\small
    \begin{tabular}{lccc}
      \toprule
      Especificación & $\beta$ (salto) & (SE) & $N$ \\
      \midrule
      Mayor único (todo)              & $-0.050^{**}$ & $(0.024)$ & 674 \\
      \quad celda \emph{treated-post} & $-0.048^{*}$  & $(0.025)$ & 639 \\
      Partido puro (sin coalición)    & $-0.032$      & $(0.030)$ & 451 \\
      \bottomrule
    \end{tabular}
  \end{table}

  \vspace{0.3em}
  {\footnotesize \textcolor{ITAMgray}{%
    \textit{No} es el efecto causal de la reforma: el doble diferencial es
    inestimable por falta de observaciones pre-reforma. $p$ de \emph{treated-post}
    por \emph{wild cluster bootstrap}. $^{**}\,p<0.05$, $^{*}\,p<0.10$.}}

\end{frame}

% ── A3 — nota sobre inferencia ────────────────────────────────────────────────
\begin{frame}[noframenumbering]{A3. Nota sobre inferencia}

  \begin{itemize}
    \item El tratamiento se decide a nivel estatal: la variación efectiva vive
          en $\approx 29$ conglomerados
    \item Con pocos conglomerados, los SE asintóticos son poco confiables; se
          acompañan con \emph{wild cluster bootstrap} (Rademacher, $B = 9{,}999$)
    \item La densidad de Cattaneo, Jansson y Ma (2020) se corre por partido sobre
          el margen firmado del partido focal, no sobre el panel apilado
          (simétrico por construcción)
    \item El balance trata covariables \emph{predeterminadas} (voto incumbente
          rezagado, estructura municipal de 2010) como desenlaces ficticios
  \end{itemize}

\end{frame}

\end{document}
